\subsection{弗雷歇空间之间的弗雷德霍姆映射与里斯映射}

命题 11. --- 设 E 与 F 为弗雷歇空间。为使从 E 到 F 中的连续线性映射 u 是弗雷德霍姆映射，必须且只需它的核与余核都是有限维的。

这由 III, p. 42 的 Prop. 2 的条件 (ii) 所给出的弗雷德霍姆映射的特征刻画以及下面的引理得出：

引理 6. --- 设 E 与 F 为弗雷歇空间，且 \(u \in \mathcal{L}(\mathrm{E};\mathrm{F})\)。假设 \(u\) 的像在 F 中有有限余维。那么 \(u\) 的像在 F 中是闭的，且 \(u\) 是严格态射。

设 G 是 F 的一个向量子空间，它是 \(\operatorname{Im}(u)\) 的补；子空间 G 是闭的（EVT, I, p. 14, cor. 1），因而商空间 F/G 是弗雷歇空间（TG, IX, p. 25, prop. 4）。设 \(\pi: F \to F/G\) 是典范满射。映射 \(\pi \circ u\) 是满射，因而是严格的（EVT, I, p. 19, cor. 3）。断言随即由 III, p. 40 的 cor. 2 得出。

命题 12. --- 设 E 与 F 为弗雷歇空间。给它们的对偶 E' 与 F' 赋以弱拓扑、紧收敛拓扑或有界收敛拓扑。设 \(u \in \mathcal{L}(\mathrm{E};\mathrm{F})\)。为使 u 是从 E 到 F 的弗雷德霍姆映射，必须且只需 \({}^{t}u\) 是从 F' 到 E' 的弗雷德霍姆映射。

该条件是必要的（III, p. 43, 第 3 目）。我们来证明它是充分的。如果 \(^{t}u\) 是弗雷德霍姆映射，那么它是严格态射（III, p. 42, prop. 2），因而 u 是有闭像的严格态射（EVT, IV, p. 28, th. 1 及 p. 29, cor. 3）。考虑到典范同构 Coker(\(^{t}u\)) → Ker(u)' 与 Ker(\(^{t}u\)) → Coker(u)'（EVT, IV, p. 27, prop. 2），向量空间 Ker(u) 与 Coker(u) 是有限维的，且 u 是弗雷德霍姆算子（III, p. 42 的 prop. 2）。

命题 13. --- 设 E 为弗雷歇空间，且 \(u \in \mathcal{L}(\mathrm{E})\)。

\begin{enumerate}
\def\labelenumi{\alph{enumi})}
\item
  自同态 u 是里斯自同态，当且仅当它的幂零空间 N 是有限维的且它的余幂零空间 I 有有限余维；
\item
 给 E' 赋以有界收敛拓扑、紧收敛拓扑或弱拓扑。如果 \(^{t}\) u 是 E' 的里斯自同态，那么 u 是 E 的里斯自同态。
\end{enumerate}
a) 假设幂零空间 N 是有限维的且余幂零空间 I 有有限余维。由于 \(\operatorname{Ker}(u) \subset \operatorname{N}\) 与 \(\operatorname{I} \subset \operatorname{Im}(u)\)，u 的核是有限维的，且 u 的像有有限余维。由定义与命题 11，映射 u 是里斯自同态。逆命题由定义得出。

b) 由命题 12，该假设蕴含 \(u\) 是弗雷德霍姆自同态，其指标为 \(\text{ind}(^t u) = -\text{ind}(u) = 0\)（\(n^o 3\)，公式 (4)）。设 \(n \in \mathbf{N}\)。由于 \(u^n\) 的像是闭的且以 \(^t u^n\) 的核为其正交补（EVT, IV, p. 27, prop. 2），序列 \((\text{Im}(u^n))\) 是平稳的，因而 \(u\) 是里斯自同态。


